
\documentclass{article}

\usepackage[utf8]{inputenc} % solo si usas pdflatex
\usepackage[T1]{fontenc}    % solo si usas pdflatex

\usepackage{amsmath,amssymb}
\usepackage{microtype}
\usepackage{cancel}         % si lo usas
\usepackage{fontawesome}    % si lo usas
% \usepackage{xypic}        % solo si lo usas
% \usepackage{multicol}     % normalmente no; usar columns de beamer
% \usepackage{amsthm}       % solo si defines teoremas propios
% \usepackage{scrextend}    % solo si lo usas

\begin{document}

\section*{Ejercicios prácticos de criptografía asimétrica}

\section{Cifrado RSA paso a paso}

\textbf{Objetivo:} Comprender el funcionamiento del algoritmo RSA aplicando cada uno de sus pasos de forma manual utilizando números pequeños.

\bigskip

Se desea cifrar y descifrar un mensaje sencillo utilizando el algoritmo de cifrado \textbf{RSA}.

\begin{enumerate}
    \item Elige los siguientes números primos:
    \[
    p = 11 \qquad q = 17
    \]

    \item Calcula:
    \begin{itemize}
        \item \( n = p \cdot q \)
        \item \( \varphi(n) = (p-1)(q-1) \)
    \end{itemize}

    \item Elige un exponente público \( e \) tal que:
    \begin{itemize}
        \item \( 1 < e < \varphi(n) \)
        \item \( \gcd(e, \varphi(n)) = 1 \)
    \end{itemize}

    \item Calcula el exponente privado \( d \) que satisfaga:
    \[
    e \cdot d \equiv 1 \pmod{\varphi(n)}
    \]

    \item Se desea cifrar el siguiente mensaje:
    \[
    m = 88
    \]
    Comprueba que se cumple \( m < n \).

    \item Calcula el mensaje cifrado utilizando la clave pública:
    \[
    c = m^e \bmod n
    \]

    \item Descifra el mensaje utilizando la clave privada:
    \[
    m = c^d \bmod n
    \]

    \item Comprueba que el mensaje descifrado coincide con el mensaje original.
\end{enumerate}

\section{Cifrado RSA con OpenSSL en WSL}

\textbf{Objetivo:} Aplicar criptografía asimétrica utilizando OpenSSL desde la línea de comandos, sin necesidad de instalar aplicaciones adicionales fuera del Subsistema de Windows para Linux (WSL).

\bigskip

Se desea cifrar y descifrar un mensaje utilizando \textbf{RSA con OpenSSL}, trabajando exclusivamente desde una terminal de WSL.

\begin{enumerate}
    \item Abre una terminal del Subsistema de Windows para Linux (WSL).

    \item Genera un par de claves RSA de 2048 bits:
    \begin{itemize}
        \item Una clave privada.
        \item A partir de ella, la clave pública.
    \end{itemize}

    \item Crea un archivo de texto llamado \texttt{mensaje.txt} que contenga un mensaje corto (por ejemplo, \emph{``Mensaje secreto''}).

    \item Cifra el archivo \texttt{mensaje.txt} utilizando:
    \begin{itemize}
        \item La clave pública generada.
        \item El algoritmo RSA mediante OpenSSL.
    \end{itemize}

    \item Comprueba que el contenido del archivo cifrado no es legible visualmente.

    \item Descifra el archivo cifrado utilizando la clave privada.

    \item Verifica que el contenido descifrado coincide exactamente con el mensaje original.

    \item (Opcional) Elimina la clave privada y reflexiona:
    \begin{itemize}
        \item ¿Sería posible recuperar el mensaje cifrado sin la clave privada?
        \item Justifica tu respuesta.
    \end{itemize}
\end{enumerate}

\bigskip

\textbf{Entrega:}
\begin{itemize}
    \item Comandos utilizados.
    \item Salida relevante obtenida en la terminal.
    \item Breve explicación del proceso seguido.
\end{itemize}


\end{document}